
\medskip
\phantomsection\label{lem:sga3-VIB-10.3}%
\noindent\textbf{Lemma 10.3.---}
\emph{Let \(G\) be an \(S\)-group of finite presentation.
Then there are an index \(j\in I\), an \(S_j\)-group \(G_j\)
of finite presentation, and an isomorphism of \(S\)-groups}
\[
  G\xrightarrow{\ \sim\ }G_j\times_{S_j}S.
\]

Assertions \hyperref[lem:sga3-VIB-10.2]{10.2} and
\hyperref[lem:sga3-VIB-10.3]{10.3} are easy consequences of
\hyperref[par:sga3-VIB-10.0]{10.0} and
\hyperref[par:sga3-VIB-10.1]{10.1}, taking into account the
interpretation%
\footnote{\label{note:sga3-VIB-93}Editor's note: in terms of
commutative diagrams of \(S\)-morphisms.}
of the structure of an \(S\)-group given in EGA~0\(_{\mathrm{III}}\),
8.2.5 and 8.2.6.

\medskip
\phantomsection\label{lem:sga3-VIB-10.4}%
\noindent\textbf{Lemma 10.4.---}
\emph{Let \(u:G\to H\) be a morphism of \(S\)-groups between
\(S\)-groups of finite presentation. By
\hyperref[lem:sga3-VIB-10.3]{10.3} and
\hyperref[lem:sga3-VIB-10.2]{10.2}, \(u\) is obtained by base
change from a morphism \(u_j:G_j\to H_j\) between
\(S_j\)-groups of finite presentation. Then \(u\) is a central
monomorphism \textup{(respectively, an invariant monomorphism)}
if and only if there is an \(i\geq j\) such that}
\phantomsection\label{form:sga3-VIB-10.4-base-change}%
\[
  u_i=u_j\times_{S_j}S_i
\]
\emph{has the same property.}

This follows immediately from
\hyperref[par:sga3-VIB-10.0]{10.0} and
\hyperref[par:sga3-VIB-10.1]{10.1}, in view of the
characterization in \hyperref[rem:sga3-VIB-6.7]{6.7} of
central or invariant group monomorphisms.

\medskip
\phantomsection\label{cor:sga3-VIB-10.5}%
\noindent\textbf{Corollary 10.5.---}
\emph{Let \(G_j\) be an \(S_j\)-group of finite presentation.
Then \(G_j\times_{S_j}S\) is commutative if and only if there
is an \(i\geq j\) such that \(G_j\times_{S_j}S_i\) is
commutative.}

Indeed, to say that an \(S\)-group is commutative is the same
as saying that it is central when regarded as a group
subscheme of itself.

\medskip
\phantomsection\label{prop:sga3-VIB-10.6}%
\noindent\textbf{Proposition 10.6.---}
\emph{Let \(G_j\) be an \(S_j\)-group of finite presentation,
and let \(G'_j\) be a group subscheme of \(G_j\), flat and of
finite presentation over \(S_j\). Then the quotient}
\phantomsection\label{form:sga3-VIB-10.6-quotient-over-S}%
\[
  (G_j\times_{S_j}S)/(G'_j\times_{S_j}S)
\]
\emph{is representable for the \textup{(fpqc)} topology if
and only if there is an \(i\geq j\) such that}
\phantomsection\label{form:sga3-VIB-10.6-quotient-over-Si}%
\[
  (G_j\times_{S_j}S_i)/(G'_j\times_{S_j}S_i)
\]
\emph{is representable.}

This is a consequence of the following more general lemma.

\medskip
\phantomsection\label{lem:sga3-VIB-10.7}%
\noindent\textbf{Lemma 10.7.---}
\emph{Let \(X_j\) be an \(S_j\)-scheme of finite presentation,
and let \(R_j\) be an equivalence relation on \(X_j\), flat and
of finite presentation.}%
\footnote{\label{note:sga3-VIB-94}Editor's note: that is, such
that the composite
\[
  R_j\lhook\joinrel\longrightarrow X_j\times_{S_j}X_j
  \xrightarrow{\ \operatorname{pr}_1\ }X_j
\]
is flat and of finite presentation.}
\emph{Then the quotient sheaf}
\[
  (X_j\times_{S_j}S)/(R_j\times_{S_j}S)
\]
\emph{for the topology \(\mathcal T=\textup{(fppf)}\) or
\(\textup{(fpqc)}\) is representable if and only if there is
an \(i\geq j\) such that the quotient sheaf}
\[
  (X_j\times_{S_j}S_i)/(R_j\times_{S_j}S_i)
\]
\emph{for the topology \(\mathcal T\) is representable.}

In view of the statements EGA~IV\(_3\),\footnote{Source note: both the current re-edition and the official original
print IV\(_2\), but all four cited sections belong to EGA~IV\(_3\).}
8.8.2, 8.8.3,
8.10.5, and 11.2.6 recalled in
\hyperref[par:sga3-VIB-10.0]{10.0}, this lemma follows from
the next result.

\medskip
\phantomsection\label{lem:sga3-VIB-10.8}%
\noindent\textbf{Lemma 10.8.---}
\emph{Let \(\mathcal T\) be the \textup{(fppf)} or
\textup{(fpqc)} topology. Let \(X\) be an \(S\)-scheme of
finite presentation \textup{(respectively, locally of finite
presentation)}, and let \(R\) be an equivalence relation on
\(X\) defined by a monomorphism}
\[
  v:R\longrightarrow X\times_S X
\]
\emph{such that \(\operatorname{pr}_1\circ v\) is flat and of
finite presentation \textup{(respectively, flat and locally
of finite presentation)}. Then the following conditions are
equivalent:}
\begin{enumerate}
\item[\textup{(i)}]\phantomsection\label{concl:sga3-VIB-10.8-i}
  \emph{The quotient sheaf \(X/R\) for the topology
  \(\mathcal T\) is representable.}

\item[\textup{(ii)}]\phantomsection\label{concl:sga3-VIB-10.8-ii}
  \emph{There are an \(S\)-scheme \(Y\) of finite presentation
  \textup{(respectively, locally of finite presentation)} and
  a faithfully flat morphism \(p:X\to Y\) such that the
  diagram}
\end{enumerate}
\begin{center}
  \phantomsection\label{diag:sga3-VIB-10.8-D}%
  \begin{tabular}{c@{\qquad}c}
    \textup{(D)} &
    \begin{tikzcd}[row sep=large,column sep=large]
      R
        \arrow[r,"{\operatorname{pr}_1\circ v}"]
        \arrow[d,"{\operatorname{pr}_2\circ v}"'] &
      X \arrow[d,"p"] \\
      X \arrow[r,"p"'] &
      Y
    \end{tikzcd}
  \end{tabular}
\end{center}
\emph{is cartesian.}

First note that, by IV, \SGARefLink{sga3:iv:definition:3:3:2}{3.3.2} and 4.4.3, the sheaf \(X/R\) for
the topology \(\mathcal T\) is representable by \(Y\) if and
only if diagram~\hyperref[diag:sga3-VIB-10.8-D]{\textup{(D)}}
is cartesian and \(p\) is a covering for the topology
\(\mathcal T\).

Let us prove that
\hyperref[concl:sga3-VIB-10.8-i]{\textup{(i)}} implies
\hyperref[concl:sga3-VIB-10.8-ii]{\textup{(ii)}}.
Hypothesis~\hyperref[concl:sga3-VIB-10.8-i]{\textup{(i)}}
implies that diagram~\hyperref[diag:sga3-VIB-10.8-D]%
{\SGARefLink{diag:sga3-VIB-10.8-D}{\textup{(D)}}} is cartesian. Thus
\(\operatorname{pr}_1\circ v\) is obtained from \(p\) by base
change along \(p\), and \(p\) is a covering for the
\textup{(fpqc)} topology. Consequently, by \textup{(fpqc)}
descent (EGA~IV\(_2\), 2.7.1), since
\(\operatorname{pr}_1\circ v\) is faithfully flat and
\textup{(locally)} of finite presentation, the same is true
of \(p\). Then EGA~IV\(_4\), 17.7.5 shows that, since \(X\) is
\textup{(locally)} of finite presentation over \(S\), so is
\(Y\).

Let us prove that
\hyperref[concl:sga3-VIB-10.8-ii]{\textup{(ii)}} implies
\hyperref[concl:sga3-VIB-10.8-i]{\textup{(i)}}. It is enough
to show that \(p\) is a covering for the \textup{(fppf)}
topology. But \(p\) is faithfully flat by hypothesis, and it
is locally of finite presentation because \(X\) and \(Y\) are
locally of finite presentation over \(S\)
\textup{(EGA~IV\(_1\), 1.4.3\textup{(v)})}.
